\chapter{Design}
\label{chap:design}
\IMRADlabel{methods} % Mechanismus um das Strukturierungsmodell IMRaD zu kennzeichnen

This chapter presents the design of the proposed ML-based view advisor. It first states the concrete objectives and system requirements derived from the problem statement, and then describes the proposed approach: system architecture, feature engineering, cost-model development and training, and the advisor integration and evaluation plan.

\section{Objectives}
\label{sec:requirements}
The primary objective of this thesis is to build a machine learning advisor that recommends the best view strategy for an arbitrary SQL workload: (i) materialized view (MV) with full recomputation, (ii) incrementally maintained materialized view (IMV), or (iii) no view. To achieve this objective the following requirements are identified:

\begin{enumerate}
    \item \textbf{Performance metrics.} Define clear, measurable metrics that capture both \emph{view maintenance cost} and \emph{query execution time} so that different strategies can be compared on the same basis.
    \item \textbf{Dataset curation.} Produce a labeled dataset of query execution times for reads and writes across the three view configurations (MV, IMV, no view). Each dataset entry must include the query, its QEP(s), and measured execution times under controlled experimental conditions.
    \item \textbf{Feature extraction.} Provide a pipeline-based feature extractor that maps QEPs into numerical feature vectors for both read and write queries.
    \item \textbf{Write cost model.} Design and train a decision-tree-based model (LightGBM) to predict per-tuple maintenance time for write pipelines; the read-cost model is assumed available.
    \item \textbf{Advisor logic.} Combine read- and write-cost predictions to estimate total workload time and produce a lowest-cost recommendation (MV / IMV / no view).
    \item \textbf{Evaluation.} Validate the advisor on unseen workloads and measure (a) prediction accuracy of the cost models, and (b) the advisor's success rate at selecting the true optimal strategy.
\end{enumerate}

\section{Proposed Approach}
\label{sec:proposed_approach}
To address the research questions posed in this thesis, we adopt a learning-based, workload-driven approach. The design is organized into four main phases: (1) system architecture, (2) featurization of SQL queries, (3) development and training of cost models, and (4) advisor integration and evaluation.

\subsection{System Architecture}
\label{sec:system_architecture}
The proposed system comprises three primary components (see Figure~\ref{fig:system}):

\begin{itemize}
    \item \textbf{Feature extractor.} Parses a query's Query Execution Plan (QEP) and decomposes it into a set of pipelines. For each pipeline the extractor builds a numerical feature vector that summarizes operator-level properties, pipeline flow statistics, and relevant schema/table statistics.
    \item \textbf{Cost models.} Two specialized models are used:
    \begin{itemize}
        \item \emph{Read cost model:} predicts the execution time of \texttt{SELECT} queries under each view configuration. A pre-existing read-cost model will be reused for this component.
        \item \emph{Write cost model:} predicts the execution time of write queries (\texttt{INSERT}, \texttt{UPDATE}, \texttt{DELETE}) with an emphasis on per-tuple maintenance cost for view maintenance pipelines.
    \end{itemize}
    \item \textbf{View advisor.} Given a workload (a sequence or distribution of SQL statements), the advisor queries the cost models for each statement under the three view scenarios, sums predicted pipeline costs to obtain per-statement and per-workload totals, and recommends the strategy with minimal predicted total workload time.
\end{itemize}

\subsection{Feature Engineering}
\label{sec:feature_engineering}
Feature engineering is split by query type (read vs.\ write) and follows a pipeline-centric decomposition.

\paragraph{Read queries}
For reads, we adopt the feature set and extraction procedures inspired by the T3 system \cite{rieger2025t3}. Features are derived from the QEP and include:
\begin{itemize}
    \item Operator-level attributes: operator type, estimated and actual output cardinalities, estimated cost, operator-specific metrics (e.g., join method).
    \item Table/column statistics: table size, column data types, histograms or distribution summaries where available.
    \item Pipeline-level aggregates: number of operators, depth of pipeline, estimated pipeline cost.
\end{itemize}
These features are designed to capture computational complexity and to enable the read-cost model to predict whether a view will accelerate the query.

\paragraph{Write queries}
For writes we use a similar pipeline-based representation but focus on tuple-centric and maintenance-relevant features:
\begin{itemize}
    \item Tuple-flow features: fraction/percentage of tuples from the pipeline start that reach particular operators (selectivity estimates).
    \item Operator-level features: operator cardinalities, output row width, estimated intermediate sizes.
    \item View-related features: number of joins and aggregations present in the view definition, and expected delta sizes that the maintenance pipeline must process.
    \item Physical costs: index lookup counts, number of affected partitions/pages where measurable.
\end{itemize}
Additionally, execution plans will be used to extract and treat distinct pipelines separately so the model can learn pipeline-specific per-tuple costs.

\subsection{Cost Model Development and Training}
\label{sec:cost_model}
The write cost model will be implemented as a LightGBM gradient-boosted decision-tree ensemble. The model training follows these principles:

\begin{itemize}
    \item \textbf{Training objective:} predict per-tuple processing time for a pipeline. The per-tuple formulation enables scaling predictions to arbitrary input cardinalities by multiplication with predicted/actual cardinalities at inference time.
    \item \textbf{Target transformation:} to stabilize training across a wide dynamic range of per-tuple times, the target per-tuple time \(t\) will be transformed via:
    \[
        t' = -\log(t).
    \]
    \item \textbf{Training data:} the dataset is curated by running workloads (mixes of reads and writes) under the three view configurations and recording per-pipeline execution times, cardinalities, and QEP-derived features. The write dataset will include reconstructed/extracted workloads (e.g., the reconstructed Amazon Redshift workload referenced in \nameref{Initial Results}).
    \item \textbf{Evaluation metrics:} model accuracy will be measured using metrics such as Mean Absolute Percentage Error (MAPE) and Q-Error on a held-out test set, and compared against baseline heuristics (e.g., simple cardinality-based rules).
    \item \textbf{Deployment optimization:} the trained LightGBM ensemble will be compiled to native code (or efficiently embedded) to enable fast inference during advisor operation.
\end{itemize}

\subsection{Advisor Integration and Evaluation}
\label{sec:integration_evaluation}
The advisor integrates the read- and write-cost models to recommend a view strategy for a workload. The per-workload decision flow is:

\begin{enumerate}
    \item For each statement in the workload, obtain its QEP in each scenario: MV, IMV, and no view.
    \item Decompose each QEP into pipelines and extract feature vectors via the feature extractor.
    \item Predict per-pipeline per-tuple time using the appropriate cost model (read or write) and scale by predicted/actual pipeline cardinalities to obtain pipeline times.
    \item Sum pipeline times to obtain per-statement predicted times for each scenario.
    \item Aggregate per-statement times across the workload to compute total predicted workload time for MV, IMV, and no view.
    \item Recommend the strategy with the minimum predicted total workload time.
\end{enumerate}

\paragraph{Evaluation plan}
End-to-end experiments will validate the advisor by executing a suite of test workloads on the target database. For each workload:
\begin{itemize}
    \item Measure the \emph{actual} total execution time under each strategy (MV, IMV, no view).
    \item Compare the advisor's recommendation against the empirically best strategy. The primary success metric is the advisor's accuracy: the percentage of workloads for which it selected the best-performing strategy.
    \item Report secondary metrics including the advisor's predicted vs.\ actual total times, model error statistics (MAPE, Q-Error), and sensitivity analyses across workload mixes (read-heavy vs.\ write-heavy).
\end{itemize}

The next chapter describes implementation details and the experimental setup used to collect the dataset and run the evaluation.
